\documentclass[12pt]{article}

\usepackage{listings}
\usepackage{color}
\definecolor{codegreen}{rgb}{0,0.6,0}
\definecolor{codegray}{rgb}{0.5,0.5,0.5}
\definecolor{codepurple}{rgb}{0.58,0,0.82}
\definecolor{backcolour}{rgb}{0.95,0.95,0.92}
\lstdefinestyle{mystyle}{
    backgroundcolor=\color{backcolour}, % 设置背景色为浅灰色 
    commentstyle=\color{codegreen},     % 代码注释颜色为绿色
    keywordstyle=\color{magenta},       % 代码关键字为品红色
    numberstyle=\tiny\color{codegray},  % 行号字号tiny，颜色浅灰 
    stringstyle=\color{codepurple},     % 字符串为紫色
    basicstyle=\ttfamily\footnotesize,  % 等宽字体
    breakatwhitespace=false,            % 换行不局限于空格
    breaklines=true,                    % 自动换行                
    captionpos=b,                       % 代码块标题位置在底部    
    keepspaces=true,                    % 保留代码中的空格
    numbers=left,                       % 行号显示在左侧
    numbersep=5pt,                      % 行号与代码间距5pt
    showspaces=false,                   % 不显示空格符号
    showstringspaces=false,             % 不显示字符串内部的空格符号
    showtabs=false,                     % 不显示制表符          
    tabsize=2                           % 制表符为2个空格
}
\lstset{style=mystyle}

\usepackage{algorithm} 
\usepackage{algorithmic} 
% \usepackage[linesnumbered, boxed]{algorithm2e} % 带行号和框线
\usepackage{amsmath, amsfonts}


\begin{document}
\section{The First Section}
% 使用 verbatim 环境显示代码
Python code example:
\begin{verbatim}
    import numpy as np
    x = np.random.rand(4)
    print(x)
\end{verbatim}

\section{The Second Section}
% 使用 listings 环境显示代码
Python code example:
\begin{lstlisting}[language=Python, caption=Example Python Code]
    import numpy as np
    x = np.random.rand(4)
    print(x)
\end{lstlisting}

\section{Pseudocode for the Algorithm}
%% 使用 algorithm 和 algorithmic 显示算法
\begin{algorithm}
\renewcommand{\algorithmicrequire}{\textbf{Input:}} % Require -> Input
\renewcommand{\algorithmicensure}{\textbf{Output:}} % Ensure -> Output
\caption{Inner product of vectors}
\begin{algorithmic}[1]
    \REQUIRE $\boldsymbol{x}, \boldsymbol{y} \in \mathbb{R}^n$
    \ENSURE $c$
    \STATE $c=0$
    \FOR{$i=1$ to $n$}
        \STATE $c=c+x_i y_i$
    \ENDFOR
\end{algorithmic}
\end{algorithm}

%% 使用 algorithm2e 显示算法
% \begin{algorithm}
%     \SetKwInOut{Input}{Input}
%     \SetKwInOut{Output}{Output}
%     \caption{Inner product of vectors}
%     \Input{$\boldsymbol{x}, \boldsymbol{y} \in \mathbb{R}^n$}
%     \Output{$c$}
%     $c=0$\;
%     \For{$i=1$ to $n$}{
%         $c=c+x_i y_i$\;
%     }
% \end{algorithm}


\end{document}